\documentclass{subfiles}
\begin{document}\label{sec:burrowsWheeler}
    Introduced in \cite{burrowsWheeler1994}, 
        \Gls{bwt} is a technique that allows a better compression.
        The transformation is simple, as described in the following.
        Let \(w\) be the text we want to transmit; 
        as first step we append a end of string symbol to it, 
        let it be \(\$\), we then compute all cyclic rotations of \(w\),
        and sort all cyclic rotations lexicographically.

    For the sake of simplicity, assume the cyclic rotations to be displayed in a matrix,
        see for instance \Cref{tab:exampleBWT}. 
        Then, output of the BWT is the last column of said matrix and the row index
        \(I\) of the original text.

    \begin{example*}
        Let \(w = mississippi\). Computing all \(w\)'s cyclic rotations we get:
        \subfile{../../TikZ/Figure - Example of the BWT}

    \noindent Sorting them lexicographically, and displaying them in a matrix we get the following.
        \subfile{../../TikZ/Table - Example of the BWT}
        
    \noindent Therefore, the output is \(\bwt = <ipssm\$ pissii, 6>\).            
    \end{example*}
    
    The BWT satisfies some interesting properties, 
        that are fundamental to compute the inverse transformation.
        These, listed below, allow us to compute two distinct mappings, namely,
        the FL and the LF mapping.
        \begin{enumerate}
            \item For every \(i \ne I, L[i]\) is followed by \(F[i]\).
            \item The first character of \(w\) is \(F[I]\).
            \item \label{bwt:3} The \(i\)-th occurence of character \(x\) in \(L\)
                corresponds to the \(i\)-th occurence of \(x\) in \(F\).
        \end{enumerate}

    Assume \(<L, I>\) to be the output of the BWT of some string. 
        From it, computing \(F\) is trivial -- 
        simply sort lexicographically the symbols in \(L\).
        By property \ref{bwt:3} we can define a mapping from \(F\) to \(L\),
        obtaining this way a permutation, call it \(\tau\).
        Then, it's easily seen that the \(i\)-th symbol of \(w\) is given by \(F[\tau^{i}[I]]\).
        The \(LF\)-mapping is obtained similarly, but the string obtained is in reverse.

    \begin{example*}
        Consider \(<h\$ lpae, 1>\).
        By sorting \(L\), we get: \(F = [\$, a, e, h, l, p]\), from which
        \[
            \tau = \begin{pmatrix}
                0 & 1 & 2 & 3 & 4 & 5 \\ 
                1 & 4 & 5 & 0 & 2 & 3 \\ 
            \end{pmatrix}.
        \]
        Retriveing \(w\), we have: \(w[0] = F[\tau^{0}[I]] = F[I] = a, 
        w[1] = F[\tau[I]] = F[4] = l, \dots, w[5] = F[\tau^{5}[I]] = F[0] = \$\).
        In conclusion, \(w = aleph\$\) from which we remove the end of string symbol \(\$\).
    \end{example*}

    \begin{remark*}
        It is easily seen that BWT's efficiency strictly depends on the sorting of 
            the cyclic rotations. It is well known that, by using a standard sorting 
            algorithms the best we can achive is a \(\bigO{n \log n}\) time complexity,
            which is obviously too much for a pre-processing step.

        It can be proved that if we focus on sorting the suffixes of the string,
            we have effectively sorted all cyclic rotations. How should we proceed then?
            In many applications we achive this by computing the \gls{sa}; 
            that is, a sorted array of the positions of the suffixes.
            In \Cref{sec:DC3} we show how to efficiently compute the SA,
            while in the next section we focus on analyzing some theoretical aspects of BWT.
    \end{remark*}

    \subsection{Analysis of the BWT}
    \subfile{../sub/5.1 - Analysis of the BWT}

    \subsection[The DC3 algorithm]{Efficient computation of the SA}
    \subfile{../sub/5.2 - DC3 algorithm}
\end{document}
